%% pc-spectre-signal — du signal temporel (à gauche) au spectre d'amplitude (à droite).
%% Exemple du cours : s(t) = 2 + 1 sin(2 pi f1 t) + 0,5 sin(2 pi (2 f1) t + pi/2),
%% soit 2 + sin(2 pi t/T) + 0,5 cos(4 pi t/T) : l'abscisse est graduée en périodes T = 1/f1.
%% Les trois composantes sont en pointillés (vert : continue, bleu : fondamental,
%% rouge : harmonique de rang 2) ; leur somme est la courbe noire épaisse.
%% La composante continue est tracée explicitement comme une raie à 0 Hz sur le spectre.
\documentclass[tikz,border=4pt]{standalone}
\usepackage{liaisons}
\begin{document}
\begin{tikzpicture}[y=0.62cm,
  axe/.style={-{Stealth[length=5pt]}, line width=0.6pt},
  grille/.style={line width=0.3pt, black!15},
  comp/.style={line width=0.9pt, dash pattern=on 2.2pt off 1.8pt},
  somme/.style={black, line width=1.3pt, line cap=round, line join=round},
  grad/.style={font=\scriptsize, inner sep=2pt},
  eti/.style={font=\scriptsize, inner sep=2pt}]

%% =================== repère temporel ==========================================
\begin{scope}[x=1.5cm]
  \node[font=\small\bfseries, anchor=south] at (1.1,4.55) {Signal dans le temps};
  %% quadrillage
  \foreach \g in {0,0.25,...,2.25} { \draw[grille] (\g,-1.5) -- (\g,3.5); }
  \foreach \g in {-1,-0.5,...,3.5} { \draw[grille] (0,\g) -- (2.25,\g); }
  %% axes
  \draw[axe] (-0.12,0) -- (2.45,0);
  \node[grad, anchor=east] at (2.45,-0.75) {$t$ (s)};
  \draw[axe] (0,-1.6) -- (0,3.75);
  \node[grad, anchor=south east] at (-0.02,3.55) {$u$ (V)};
  \foreach \v in {-1,1,2,3} {
    \draw[line width=0.5pt] (-0.03,\v) -- (0.03,\v);
    \node[grad, anchor=east] at (-0.05,\v) {\v}; }
  \node[grad, anchor=north east] at (-0.03,0) {0};

  %% composantes
  \draw[comp, solideE] (0,2) -- (2.25,2);
  \draw[comp, solideB] plot[domain=0:2.25, samples=90, smooth] (\x,{sin(360*\x)});
  \draw[comp, solideA] plot[domain=0:2.25, samples=90, smooth] (\x,{0.5*cos(720*\x)});
  %% somme des trois
  \draw[somme] plot[domain=0:2.25, samples=140, smooth] (\x,{2+sin(360*\x)+0.5*cos(720*\x)});
  \node[eti, text=black, anchor=west] at (1.62,3.30) {$s(t)$};
  \node[eti, text=solideE, anchor=west] at (2.28,2) {$A_0$};
  \node[eti, text=solideB, anchor=west] at (2.28,1) {$A_1$};
  \node[eti, text=solideA, anchor=west] at (2.28,0.5) {$A_2$};

  %% période
  \draw[line width=0.4pt, black!55, dash pattern=on 1.6pt off 1.6pt] (0.25,2.5) -- (0.25,4.0);
  \draw[line width=0.4pt, black!55, dash pattern=on 1.6pt off 1.6pt] (1.25,2.5) -- (1.25,4.0);
  \draw[{Stealth[length=4pt]}-{Stealth[length=4pt]}, line width=0.6pt] (0.25,3.85) -- (1.25,3.85);
  \node[grad, anchor=south, inner sep=1pt] at (0.75,3.90) {$T = 1/f_1$};
\end{scope}

%% =================== flèche de passage ========================================
\begin{scope}[x=1cm]
  \draw[-{Stealth[length=8pt]}, line width=3pt, solideF] (4.20,1.6) -- (5.30,1.6);
  \node[eti, text=solideF, anchor=south, align=center] at (4.60,2.15)
    {décomposition\\de Fourier};
\end{scope}

%% =================== spectre d'amplitude ======================================
\begin{scope}[shift={(5.55cm,0)}, x=0.95cm]
  \node[font=\small\bfseries, anchor=south] at (1.4,4.55) {Spectre d'amplitude};
  \draw[axe] (-0.15,0) -- (3.30,0) node[anchor=west, inner sep=3pt, font=\small] {$f$ (Hz)};
  \draw[axe] (0,-0.25) -- (0,3.75);
  \node[grad, anchor=south east] at (-0.02,3.55) {$u$ (V)};
  \foreach \v in {1,2} {
    \draw[line width=0.5pt] (-0.05,\v) -- (0.05,\v);
    \node[grad, anchor=east] at (-0.07,\v) {\v}; }
  \node[grad, anchor=north east] at (-0.03,0) {0};
  \foreach \f/\lab in {1/{$f_1$}, 2/{$f_2 = 2f_1$}} {
    \draw[line width=0.5pt] (\f,-0.12) -- (\f,0.12);
    \node[grad, anchor=north] at (\f,-0.14) {\lab}; }

  %% les trois raies
  \draw[solideE, line width=2pt] (0,0) -- (0,2);
  \draw[solideB, line width=2pt] (1,0) -- (1,1);
  \draw[solideA, line width=2pt] (2,0) -- (2,0.5);
  \foreach \f/\h/\co in {0/2/solideE, 1/1/solideB, 2/0.5/solideA} {
    \fill[\co] (\f,\h) circle (0.05cm); }

  \node[eti, text=solideE, anchor=west, align=left] at (0.12,2.55)
    {composante continue (0 Hz)\\$A_0 = 2$ V};
  \node[eti, text=solideB, anchor=west, align=left] at (1.12,1.65)
    {fondamental\\$A_1 = 1$ V};
  \node[eti, text=solideA, anchor=west, align=left] at (2.12,0.85)
    {harmonique de rang 2\\$A_2 = 0{,}5$ V};
\end{scope}
\end{tikzpicture}
\end{document}
